\setchoices{label=(\Alph*)}

\begin{question}[index=1]
设 $f(x)$ 为连续函数，且 $\displaystyle\lim_{x\to1}\frac{f(2x-1)}{x-1}=2$，则 $\displaystyle\lim_{x\to0}\frac{f[(1+\sin x)^2]-f(1+\sin x)}x=$
\begin{choices}[columns=4]
\item $0$
\item $2$
\item $1$
\item $-1$
\end{choices}
\end{question}
\QuestionBreak

\begin{question}[index=2]
下列积分收敛的是
\begin{choices}[columns=2]
\item $\displaystyle\int_0^{+\infty}\frac{\eu^{1/x}}{x^4}\dd x$
\item $\displaystyle\int_0^{\pi/2}\frac{1-\cos x}{x^3}\dd x$
\item $\displaystyle\int_{-\infty}^0\frac{\eu^{1/x}}{x^2+\sin^2x}\dd x$
\item $\displaystyle\int_0^{+\infty}\frac{\eu^{1/x}}{x^2+\sin^2x}\dd x$
\end{choices}
\end{question}
\QuestionBreak

\begin{question}[index=3]
设 $D=\{(x,y)\mid(x-1)^2+(y-1)^2\le1\}$，$\displaystyle I_k=\iint_D(x+y)^k\dd x\dd y$（$k=1,2,3$），则正确的是
\begin{choices}[columns=4]
\item $I_1<I_2<I_3$
\item $I_3<I_2<I_1$
\item $I_1<I_3<I_2$
\item $I_3<I_1<I_2$
\end{choices}
\end{question}
\QuestionBreak

\begin{question}[index=4]
设数列 $\{a_n\}$ 满足 $a_0=1$，$\displaystyle\sum_{n=0}^\infty2a_nx^{n+1}+\sum_{n=0}^\infty(n+1)a_{n+1}x^n=0$，则 $\displaystyle\sum_{n=0}^\infty a_nx^n=$
\begin{choices}[columns=4]
\item $\eu^x$
\item $\eu^{-x}$
\item $\eu^{-x^2}$
\item $\eu^{x^2}$
\end{choices}
\end{question}
\QuestionBreak

\begin{question}[index=5]
设 $\vect A$ 为 $m\times n$ 矩阵，$\vect B$ 为 $n\times s$ 矩阵，则错误的是
\begin{choices}[columns=1]
\item 若方程组 $\vect A\vect B\vect X=\vect0$ 与 $\vect B\vect X=\vect0$ 同解，则 $r(\vect A\vect B)=r(\vect B)$
\item 若 $r(\vect A\vect B)=r(\vect B)$，则方程组 $\vect A\vect B\vect X=\vect0$ 与 $\vect B\vect X=\vect0$ 同解
\item 若 $r(\vect A)=n$，则方程组 $\vect A\vect B\vect X=\vect0$ 与 $\vect B\vect X=\vect0$ 同解
\item 若方程组 $\vect A\vect B\vect X=\vect0$ 与 $\vect B\vect X=\vect0$ 同解，则 $r(\vect A)=n$
\end{choices}
\end{question}
\QuestionBreak

\begin{question}[index=6]
设 $\vect A$ 是 $3\times4$ 实矩阵，$\vect E_3$ 与 $\vect E_4$ 分别为 $3$ 阶与 $4$ 阶单位矩阵，则 $r(\vect E_4-\vect A^{\mathrm T}\vect A)-r(\vect E_3-\vect A\vect A^{\mathrm T})=$
\begin{choices}[columns=4]
\item $0$
\item $1$
\item $2$
\item $3$
\end{choices}
\end{question}
\QuestionBreak

\begin{question}[index=7]
实矩阵 $\vect A=\mat{1&0&0&0\\a&1&0&0\\1&b&2&0\\3&2&c&2}$ 与对角矩阵相似的充分必要条件为
\begin{choices}[columns=4]
\item $a-b=0$
\item $a-c=0$
\item $a=c=0$
\item $a=b=c=0$
\end{choices}
\end{question}
\QuestionBreak

\begin{question}[index=8]
设平面区域 $G=\{(x,y)\mid0\le x\le1,0\le y\le1\}$，现向 $G$ 上均匀投掷随机点，已知随机点的落点 $(X,Y)$ 在 $G$ 上分布均匀，记 $A=\{X\le\dfrac12\}$，$B=\{Y\le\dfrac12\}$，$C=\{(X-\dfrac12)(Y-\dfrac12)<0\}$，则
\begin{choices}[columns=2]
\item $A,B,C$ 相互独立
\item $AB$ 与 $C$ 相互独立
\item $A$ 与 $BC$ 相互独立
\item $A,B,C$ 两两独立但不相互独立
\end{choices}
\end{question}
\QuestionBreak

\begin{question}[index=9]
设随机变量 $X$ 在 $(0,2)$ 内服从均匀分布，在 $X=x$（$0<x<2$）的条件下，$Y$ 在区间 $(0,x)$ 内服从均匀分布，则 $\operatorname{Cov}(X,Y)=$
\begin{choices}[columns=4]
\item $1$
\item $\dfrac12$
\item $\dfrac14$
\item $\dfrac16$
\end{choices}
\end{question}
\QuestionBreak

\begin{question}[index=10]
设总体 $X$ 在区间 $[\theta,\theta+1]$ 上服从均匀分布，$\theta$（$-\infty<\theta<+\infty$）是未知参数，$X_1,X_2,\ldots,X_n$（$n>1$）为来自总体 $X$ 的简单随机样本．记 $X_{(1)}=\min(X_1,X_2,\ldots,X_n)$，$X_{(n)}=\max(X_1,X_2,\ldots,X_n)$，则不是 $\theta$ 的最大似然估计量的是
\begin{choices}[columns=2]
\item $\widehat\theta_1=X_{(1)}$
\item $\widehat\theta_2=X_{(n)}-1$
\item $\widehat\theta_3=\dfrac12(X_{(1)}+X_{(n)})-\dfrac12$
\item $\widehat\theta_4=X_{(1)}+1$
\end{choices}
\end{question}
